% Options for packages loaded elsewhere
\PassOptionsToPackage{unicode}{hyperref}
\PassOptionsToPackage{hyphens}{url}
\documentclass[
]{article}
\usepackage{xcolor}
\usepackage[margin=1in]{geometry}
\usepackage{amsmath,amssymb}
\setcounter{secnumdepth}{-\maxdimen} % remove section numbering
\usepackage{iftex}
\ifPDFTeX
  \usepackage[T1]{fontenc}
  \usepackage[utf8]{inputenc}
  \usepackage{textcomp} % provide euro and other symbols
\else % if luatex or xetex
  \usepackage{unicode-math} % this also loads fontspec
  \defaultfontfeatures{Scale=MatchLowercase}
  \defaultfontfeatures[\rmfamily]{Ligatures=TeX,Scale=1}
\fi
\usepackage{lmodern}
\ifPDFTeX\else
  % xetex/luatex font selection
\fi
% Use upquote if available, for straight quotes in verbatim environments
\IfFileExists{upquote.sty}{\usepackage{upquote}}{}
\IfFileExists{microtype.sty}{% use microtype if available
  \usepackage[]{microtype}
  \UseMicrotypeSet[protrusion]{basicmath} % disable protrusion for tt fonts
}{}
\makeatletter
\@ifundefined{KOMAClassName}{% if non-KOMA class
  \IfFileExists{parskip.sty}{%
    \usepackage{parskip}
  }{% else
    \setlength{\parindent}{0pt}
    \setlength{\parskip}{6pt plus 2pt minus 1pt}}
}{% if KOMA class
  \KOMAoptions{parskip=half}}
\makeatother
\usepackage{color}
\usepackage{fancyvrb}
\newcommand{\VerbBar}{|}
\newcommand{\VERB}{\Verb[commandchars=\\\{\}]}
\DefineVerbatimEnvironment{Highlighting}{Verbatim}{commandchars=\\\{\}}
% Add ',fontsize=\small' for more characters per line
\usepackage{framed}
\definecolor{shadecolor}{RGB}{248,248,248}
\newenvironment{Shaded}{\begin{snugshade}}{\end{snugshade}}
\newcommand{\AlertTok}[1]{\textcolor[rgb]{0.94,0.16,0.16}{#1}}
\newcommand{\AnnotationTok}[1]{\textcolor[rgb]{0.56,0.35,0.01}{\textbf{\textit{#1}}}}
\newcommand{\AttributeTok}[1]{\textcolor[rgb]{0.13,0.29,0.53}{#1}}
\newcommand{\BaseNTok}[1]{\textcolor[rgb]{0.00,0.00,0.81}{#1}}
\newcommand{\BuiltInTok}[1]{#1}
\newcommand{\CharTok}[1]{\textcolor[rgb]{0.31,0.60,0.02}{#1}}
\newcommand{\CommentTok}[1]{\textcolor[rgb]{0.56,0.35,0.01}{\textit{#1}}}
\newcommand{\CommentVarTok}[1]{\textcolor[rgb]{0.56,0.35,0.01}{\textbf{\textit{#1}}}}
\newcommand{\ConstantTok}[1]{\textcolor[rgb]{0.56,0.35,0.01}{#1}}
\newcommand{\ControlFlowTok}[1]{\textcolor[rgb]{0.13,0.29,0.53}{\textbf{#1}}}
\newcommand{\DataTypeTok}[1]{\textcolor[rgb]{0.13,0.29,0.53}{#1}}
\newcommand{\DecValTok}[1]{\textcolor[rgb]{0.00,0.00,0.81}{#1}}
\newcommand{\DocumentationTok}[1]{\textcolor[rgb]{0.56,0.35,0.01}{\textbf{\textit{#1}}}}
\newcommand{\ErrorTok}[1]{\textcolor[rgb]{0.64,0.00,0.00}{\textbf{#1}}}
\newcommand{\ExtensionTok}[1]{#1}
\newcommand{\FloatTok}[1]{\textcolor[rgb]{0.00,0.00,0.81}{#1}}
\newcommand{\FunctionTok}[1]{\textcolor[rgb]{0.13,0.29,0.53}{\textbf{#1}}}
\newcommand{\ImportTok}[1]{#1}
\newcommand{\InformationTok}[1]{\textcolor[rgb]{0.56,0.35,0.01}{\textbf{\textit{#1}}}}
\newcommand{\KeywordTok}[1]{\textcolor[rgb]{0.13,0.29,0.53}{\textbf{#1}}}
\newcommand{\NormalTok}[1]{#1}
\newcommand{\OperatorTok}[1]{\textcolor[rgb]{0.81,0.36,0.00}{\textbf{#1}}}
\newcommand{\OtherTok}[1]{\textcolor[rgb]{0.56,0.35,0.01}{#1}}
\newcommand{\PreprocessorTok}[1]{\textcolor[rgb]{0.56,0.35,0.01}{\textit{#1}}}
\newcommand{\RegionMarkerTok}[1]{#1}
\newcommand{\SpecialCharTok}[1]{\textcolor[rgb]{0.81,0.36,0.00}{\textbf{#1}}}
\newcommand{\SpecialStringTok}[1]{\textcolor[rgb]{0.31,0.60,0.02}{#1}}
\newcommand{\StringTok}[1]{\textcolor[rgb]{0.31,0.60,0.02}{#1}}
\newcommand{\VariableTok}[1]{\textcolor[rgb]{0.00,0.00,0.00}{#1}}
\newcommand{\VerbatimStringTok}[1]{\textcolor[rgb]{0.31,0.60,0.02}{#1}}
\newcommand{\WarningTok}[1]{\textcolor[rgb]{0.56,0.35,0.01}{\textbf{\textit{#1}}}}
\usepackage{graphicx}
\makeatletter
\newsavebox\pandoc@box
\newcommand*\pandocbounded[1]{% scales image to fit in text height/width
  \sbox\pandoc@box{#1}%
  \Gscale@div\@tempa{\textheight}{\dimexpr\ht\pandoc@box+\dp\pandoc@box\relax}%
  \Gscale@div\@tempb{\linewidth}{\wd\pandoc@box}%
  \ifdim\@tempb\p@<\@tempa\p@\let\@tempa\@tempb\fi% select the smaller of both
  \ifdim\@tempa\p@<\p@\scalebox{\@tempa}{\usebox\pandoc@box}%
  \else\usebox{\pandoc@box}%
  \fi%
}
% Set default figure placement to htbp
\def\fps@figure{htbp}
\makeatother
\setlength{\emergencystretch}{3em} % prevent overfull lines
\providecommand{\tightlist}{%
  \setlength{\itemsep}{0pt}\setlength{\parskip}{0pt}}
\usepackage{bookmark}
\IfFileExists{xurl.sty}{\usepackage{xurl}}{} % add URL line breaks if available
\urlstyle{same}
\hypersetup{
  pdftitle={The Posterior Predictive Distribution},
  hidelinks,
  pdfcreator={LaTeX via pandoc}}

\title{The Posterior Predictive Distribution}
\author{}
\date{\vspace{-2.5em}2026-04-17}

\begin{document}
\maketitle

\subsection{Introduction}\label{introduction}

The posterior predictive distribution (PPD) is the Bayesian answer to
the question ``what range of values can a future observation plausibly
take?'' It accounts for two sources of uncertainty simultaneously: the
sampling variability that would persist even if the parameters were
known, and the residual parameter uncertainty quantified by the
posterior. A plug-in prediction that uses only the posterior mean throws
away the second source and produces intervals that are too narrow --- a
frequent and consequential mistake in applied work.

\subsection{Prerequisites}\label{prerequisites}

A working understanding of posterior distributions, the difference
between parameter inference and observation prediction, and Monte Carlo
sampling from MCMC draws.

\subsection{Theory}\label{theory}

For data \(y\) and a prospective new observation \(\tilde y\), the
posterior predictive distribution is

\[p(\tilde y \mid y) = \int p(\tilde y \mid \theta) p(\theta \mid y) \, \mathrm d \theta.\]

The Monte Carlo approximation is direct: for each posterior draw
\(\theta^{(s)}\), sample
\(\tilde y^{(s)} \sim p(\tilde y \mid \theta^{(s)})\). The collection
\(\{\tilde y^{(s)}\}_{s=1}^S\) is itself a sample from the PPD, and any
quantile, interval, or expectation of interest comes from it. Because
the integral over \(\theta\) is automatic in this formulation, the
resulting prediction interval is wider than the corresponding interval
based on a plug-in \(\hat\theta\), and that extra width is the
propagated parameter uncertainty.

\subsection{Assumptions}\label{assumptions}

Valid posterior draws (converged chains, adequate ESS) and the assumed
data-generating model \(p(\tilde y \mid \theta)\) for the new
observation. If the new observation is from a different population or
covariate range, the PPD inherits the model's extrapolation behaviour
and may be unrealistic.

\subsection{R Implementation}\label{r-implementation}

\begin{Shaded}
\begin{Highlighting}[]
\FunctionTok{library}\NormalTok{(brms)}

\FunctionTok{set.seed}\NormalTok{(}\DecValTok{2026}\NormalTok{)}
\NormalTok{n }\OtherTok{\textless{}{-}} \DecValTok{100}
\NormalTok{d }\OtherTok{\textless{}{-}} \FunctionTok{data.frame}\NormalTok{(}\AttributeTok{x =} \FunctionTok{rnorm}\NormalTok{(n), }\AttributeTok{y =} \ConstantTok{NA}\NormalTok{)}
\NormalTok{d}\SpecialCharTok{$}\NormalTok{y }\OtherTok{\textless{}{-}} \DecValTok{2} \SpecialCharTok{+} \FloatTok{1.5} \SpecialCharTok{*}\NormalTok{ d}\SpecialCharTok{$}\NormalTok{x }\SpecialCharTok{+} \FunctionTok{rnorm}\NormalTok{(n)}

\NormalTok{fit }\OtherTok{\textless{}{-}} \FunctionTok{brm}\NormalTok{(y }\SpecialCharTok{\textasciitilde{}}\NormalTok{ x, }\AttributeTok{data =}\NormalTok{ d, }\AttributeTok{chains =} \DecValTok{2}\NormalTok{, }\AttributeTok{iter =} \DecValTok{2000}\NormalTok{, }\AttributeTok{refresh =} \DecValTok{0}\NormalTok{)}

\CommentTok{\# Posterior predictive samples}
\NormalTok{pp }\OtherTok{\textless{}{-}} \FunctionTok{posterior\_predict}\NormalTok{(fit, }\AttributeTok{newdata =} \FunctionTok{data.frame}\NormalTok{(}\AttributeTok{x =} \DecValTok{0}\SpecialCharTok{:}\DecValTok{3}\NormalTok{))}
\FunctionTok{apply}\NormalTok{(pp, }\DecValTok{2}\NormalTok{, quantile, }\FunctionTok{c}\NormalTok{(}\FloatTok{0.025}\NormalTok{, }\FloatTok{0.5}\NormalTok{, }\FloatTok{0.975}\NormalTok{))}
\end{Highlighting}
\end{Shaded}

\subsection{Output \& Results}\label{output-results}

\texttt{posterior\_predict()} returns a matrix with one row per
posterior draw and one column per \texttt{newdata} row. Summarising
columnwise gives the posterior predictive median and 95 \% interval at
each new \(x\). The resulting interval combines the residual standard
deviation with the parameter uncertainty inherited from the posterior of
\(\beta_0\) and \(\beta_1\).

\subsection{Interpretation}\label{interpretation}

A reporting sentence: ``At \(x = 2\) the posterior predictive 95 \%
interval for a new observation is 3.1 to 6.8, with median 5.0; the
corresponding interval for the conditional mean is 4.7 to 5.3,
illustrating that residual variability dominates the predictive
uncertainty in this model.'' Distinguishing the predictive interval from
the credible interval on the mean is essential --- they answer different
questions and the former is always wider.

\subsection{Practical Tips}\label{practical-tips}

\begin{itemize}
\tightlist
\item
  Use \texttt{posterior\_epred()} for the conditional-mean credible
  interval (no residual variability), \texttt{posterior\_predict()} for
  the full predictive distribution including residual variation, and
  \texttt{posterior\_linpred()} if you need the linear predictor before
  any link transformation.
\item
  Pair PPD draws with \texttt{pp\_check(fit)} to confirm that the model
  captures the marginal distribution of the data; a misspecified family
  will produce predictive intervals that are too tight or too wide where
  they matter.
\item
  For binary, count, or ordinal outcomes, the PPD is a discrete
  distribution; visualise with rootograms or bar charts rather than
  density curves.
\item
  When predicting beyond the observed predictor range, report posterior
  predictive intervals separately for in-sample and out-of-sample \(x\);
  readers should see where extrapolation begins.
\item
  Always summarise the PPD on the original scale of the outcome;
  transforming a credible interval on the linear predictor and adding
  back-of-the-envelope residual SD is rarely correct.
\end{itemize}

\subsection{R Packages Used}\label{r-packages-used}

\texttt{brms} and \texttt{rstanarm} for \texttt{posterior\_predict()},
\texttt{posterior\_epred()}, and \texttt{posterior\_linpred()};
\texttt{tidybayes::predicted\_draws()} for tidy long-format predictive
draws; \texttt{bayesplot} for \texttt{pp\_check()} and other PPD
diagnostic plots.

\subsection{For Reviewers}\label{for-reviewers}

\textbf{What to look for in a paper using this method.}

\begin{itemize}
\tightlist
\item
  \textbf{Common misapplications.}
\item
  \textbf{Diagnostics that should be reported but often aren't.}
\item
  \textbf{Red flags in tables and figures.}
\item
  \textbf{What to verify.}
\item
  \textbf{What an adequate Methods paragraph must contain.}
\end{itemize}

\subsection{See also --- labs in this
chapter}\label{see-also-labs-in-this-chapter}

\begin{itemize}
\tightlist
\item
  \href{labs/lab-bayesian-thinking-control-vs-decision-error.Rmd}{Bayesian
  thinking; α control vs decision error}
\item
  \href{labs/lab-brms-stan-loo-hierarchical-models.Rmd}{brms / Stan,
  LOO, hierarchical models}
\end{itemize}

\end{document}
